% ============================================================
% 01_avance.tex — Avance realizado esta semana
% Responsable: TODO (asignar integrante)
% ============================================================

\section{Avance Realizado}

\subsection{Lectura y comprensión del paper}

Leímos el paper completo de C-Lock. Lo que más nos costó entender al principio
fue cómo el hardware puede detener el reloj de un core sin que este pierda su
estado — el mecanismo de \textit{clock-gating} no es tan directo como parece
desde el software. Una vez que lo vimos con más detalle, identificamos los
componentes principales:

\begin{itemize}
    \item \textbf{C-LockManager}: el componente central. Tiene un árbitro que
          decide quién accede al lock, un contador global de locks activos, y
          un pool por cada core.
    \item \textbf{Pool\_i / Conflict Checker}: cada core tiene su propio pool
          con un \textit{Conflict Checker} que detecta si hay conflicto para
          un lock dado y activa el clock-gating si corresponde.
    \item \textbf{Interfaz SW/HW}: el programador usa tres macros —
          \texttt{ADD\_ITEM(id)}, \texttt{BEGIN\_C\_LOCK(id)} y
          \texttt{END\_C\_LOCK(id)} — para interactuar con el hardware.
\end{itemize}

Lo interesante es que desde el software se ve casi igual que un mutex clásico,
pero por debajo el hardware hace todo el trabajo de sincronización y además
apaga el reloj del core mientras espera, ahorrando energía.

\subsection{Organización del repositorio}

Creamos el repositorio en GitHub y lo organizamos en dos carpetas principales:

\begin{itemize}
    \item \texttt{latex/}: contiene el informe dividido en secciones separadas,
          una por tema, para que cada integrante pueda trabajar su parte sin
          generar conflictos en git.
    \item \texttt{sim/}: contiene el esqueleto de la simulación en C, con
          headers en \texttt{include/} y fuentes en \texttt{src/}.
\end{itemize}

\subsection{Esqueleto de la simulación en C}

Definimos la arquitectura del simulador y escribimos el esqueleto inicial.
Los módulos que tenemos son:

\begin{itemize}
    \item \texttt{c\_lock.h}: tipos y constantes compartidas
          (\texttt{CoreState}, modelo de energía, tamaños máximos).
    \item \texttt{c\_lock\_manager.h/.c}: gestor central — maneja la tabla de
          locks, las colas de espera y notifica a los cores cuando el lock
          queda libre.
    \item \texttt{core.h/.c}: representa un core del procesador; el diseño
          contempla usar \texttt{pthread\_cond\_wait} para emular el clock-gating
          sin busy-waiting.
    \item \texttt{benchmark.c}: estructura del escenario de prueba (N cores,
          1 lock compartido); la lógica interna está pendiente de implementar.
\end{itemize}

El esqueleto compila sin errores. La implementación de la lógica de
sincronización está en progreso y es el foco principal para la semana 6.
